%-------------------------------------------------------------------------------
\section{Restricted Code Generation}
\label{sec:codegen}
%-------------------------------------------------------------------------------

Generating checking code directly with a large language model faces two
inherent difficulties. The first is hallucination. The model may
misspell a field name as another field that happens to exist but is
semantically unrelated, or introduce a unit error in an OID. The second is
output instability. The same code-generation prompt often produces
different outputs across calls, with no guarantee that the result will compile or
be logically correct. To contain this uncertainty, we confine a rule's code body
to a finite, closed DSL-tree space $\Tv$. Each rule's code body can only be a
tree in that space, which is then rendered into Go code by a deterministic
function. The deterministic main path avoids unconstrained code emission, and
the restricted LLM fallback is bounded to the same closed space and gated by
``render and pass go build,'' so any out-of-space or non-compiling output is
rejected rather than shipped.
Figure~\ref{fig:codegen-pipeline} illustrates the whole
pipeline on a single normative rule---CA certificates must assert the keyCertSign
key-usage bit.

\begin{figure*}[tb]
  \centering
  \paperfig{figures/fig03_codegen_pipeline.png}{Worked example of restricted
    code generation. A single normative rule flows from its filled IR fields through the
    DSL tree, represented as an abstract syntax tree, and its JSON serialization to deterministically rendered Go
    checking code.}
  \caption{Worked example of restricted code generation. }
  \label{fig:codegen-pipeline}
\end{figure*}

\subsection{Atomic Templates and IR Filling}
\label{subsec:atoms}

An atomic template is the smallest unit of a certificate lint.
Most certificate lints in PKI are logical combinations of a
few atomic templates that test whether an extension is present, whether a
field equals a constant, or whether every element of a list matches a
pattern, and combining such atomic templates with conjunction, disjunction, and
negation can
characterize most static constraints. The abstract syntax of the code DSL is:
\begin{equation}\label{eq:grammar}
  \mathcal{T} \;::=\; a(\bar{v}) \;\mid\; \neg\, \mathcal{T} \;\mid\;
  \mathcal{T} \wedge \mathcal{T} \;\mid\; \mathcal{T} \vee \mathcal{T}
\end{equation}
where $a \in \mathcal{A}$ is an atomic-template predicate and $\bar{v}$ is its
argument list. $\mathcal{A}$ is a finite, closed set of atomic
templates; the experiments use 151 atomic templates in total, split by
generality into 98 GENERIC and 53 NON\_GENERIC, with grading criteria and
representative examples in Appendix~\ref{app:tv}; $\{\neg, \wedge, \vee\}$ are
the propositional-logic combinators. The code body of a lint rule is modeled as
an ordered pair $(p, q) \in \mathcal{T}_\varepsilon \times \mathcal{T}$, where
$q \in \mathcal{T}$ is the main assertion and $p \in \mathcal{T}_\varepsilon =
\mathcal{T} \cup \{\varepsilon\}$ is an optional precondition, where $\varepsilon$
marks an absent precondition. In the
equation below, $c$ is a single certificate, and $\lVert u \rVert(c) \in
\{\mathrm{true}, \mathrm{false}\}$ denotes the Boolean result of evaluating a
DSL tree or precondition $u$ on $c$; atomic templates evaluate by their semantics, and
$\neg/\wedge/\vee$ by classical propositional logic. The execution semantics
is:
\begin{equation}\label{eq:exec}
  \lVert (p, q) \rVert(c) \;=\;
  \begin{cases}
    \mathrm{NA}, & p \neq \varepsilon \,\land\, \lnot\lVert p \rVert(c) \\[2pt]
    \mathrm{Pass}, & (p = \varepsilon \,\lor\, \lVert p \rVert(c)) \,\land\,
      \lVert q \rVert(c) \\[2pt]
    \mathrm{Severity}(r), & \text{otherwise}
  \end{cases}
\end{equation}
The precondition $p$ is the normative rule's precondition. When it exists and is
false on $c$, the result is $\mathrm{NA}$, not applicable; when the
precondition is absent or holds and $q$ holds, the result is $\mathrm{Pass}$;
when the precondition holds but $q$ does not, it is a violation, returning
$\mathrm{Severity}(r)$ by obligation. Because $C_1$ of \S\ref{sec:lintability}
already excludes MAY/OPTIONAL, an obligation entering generation must be
MUST-class or SHOULD-class, so $\mathrm{Severity}(r) \in \{\mathrm{Error},
\mathrm{Warn}\}$.

\textbf{The IR fills the atomic template's parameter slots.} Selecting an atomic
template fixes only which check to use; one still needs to know what to check and
against what value, and these parameter values come precisely from the current
normative rule's IR. The IR's subject, the certificate field path given by the
field resolver such as subjectAltName.dNSName, fills the atomic
template's field slot, and the IR's constraint, a constraint value or
pattern such as a value set, a length interval, or a regex name, fills the value
slot. For example, the normative rule ``the dNSName of subjectAltName must be non-empty''
is, after mapping and filling, instantiated as
$\mathrm{FieldNonEmpty}(\mathrm{DNSNames})$.

\textbf{Mapping IR predicates to atomic templates.} The upstream
IR and the downstream DSL use two vocabularies. The IR describes normative rules
with extraction-stage predicates such as must\_be\_present,
encode\_as, and in\_range, while the DSL expresses checks with
atomic templates, and the two are not in one-to-one correspondence. To
bridge them, we maintain a many-to-many predicate-to-template map that assigns each IR
predicate a set of candidate atomic templates that can semantically carry it.
For instance, must\_be\_present may correspond to $\{\mathrm{ExtPresent},
\mathrm{FieldNonEmpty}\}$, and in\_range to $\{\mathrm{IntInRange},
\mathrm{PathLenConstraintHas}, \dots\}$---which one to take depends on the type
of the constrained field.

\subsection{Typed Vocabulary and Parameter Closure}
\label{subsec:tv}

\S\ref{subsec:atoms} described how atomic templates and IR filling place the
IR's subject and constraint into an atomic template's parameter slots---but the
IR provides only the field that a normative rule restricts and the value it
restricts that field to,
without constraining the legal type of those values. If the subject path
in the IR points to a nonexistent field, or a string is placed into a slot that
requires an integer, a structurally legal but semantically misaligned tree can
still result. This subsection closes that gap with a typed vocabulary
$\mathcal{V}$. $\mathcal{V}$ is the disjoint union of seven finite sets, frozen
at system startup, and each atomic-template parameter may only draw from one of
them---certificate field names, DN fields, OID constants, KeyUsage bits,
ExtKeyUsage bits, ASN.1 encoding types, and named regexes. The size and samples
of each class appear in Appendix~\ref{app:tv}.

Each atomic template has a signature table specifying which class, or base type
integer/Boolean/string, each of its parameters must draw from; a call to an
atomic template is legal if and only if every actual argument falls in the class its
signature prescribes. Writing $\Tv$ for the set of DSL trees all of whose
parameters are legal, we strictly restrict the range of the code generator to
it. Any tree in $\Tv$ contains no field name, OID, or regex
outside $\mathcal{V}$; an out-of-bounds tree necessarily errors at the parsing
stage, triggering repair, and never enters code generation.

\subsection{Code Generation and Mechanical Translation}
\label{subsec:mech}

A legal DSL tree $t$ must be turned into two products. One is executable Go
checking code, and the other is a code summary, the code\_summary, that
summarizes its checking logic in one sentence. Both conversions are performed by
deterministic total functions, with no LLM call. The rendering function $\rho$
maps a tree to Go code, and the mechanical translation function $\smech$ maps the
same tree to a code\_summary sentence; for the same $t$, each produces a unique
output.

$\rho$ is type-safe. It never generates an out-of-bounds field, OID, or
parameter type, so closure is preserved from the DSL-tree layer down to the
Go-code layer. Together with determinism, this lets the certificate-level
execution verification of \S\ref{subsec:exec-verify} treat $\rho(t)$ as a fixed
function of $t$, whose execution behavior does not vary with run count or model.
$\smech$, by contrast, is structure-preserving. Each atomic template and
combinator maps to a distinct phrase, so $\smech(t)$ retains the structural
information of the DSL tree in the verification chain; the construction is
detailed in Appendix~\ref{app:sigma}.

Thus $\rho$'s product leads to executable code and the certificate-level execution
verification of \S\ref{subsec:exec-verify}, while $\smech$'s product leads to
semantic-equivalence verification---comparing the code\_summary with the
normative rule text and determining whether the two are semantically equivalent.

\subsection{Tree Synthesis with a Deterministic Main Path and Restricted LLM
Fallback}
\label{subsec:synth}

The preceding discussion described how the IR filling of \S\ref{subsec:atoms} and the
parameter closure of \S\ref{subsec:tv} confine a DSL tree to $\Tv$, but two
decisions still require synthesis. The first is which atomic template carries
which IR predicate. The second is whether the IR's subject and
constraint can both be filled into the chosen atomic template.
This subsection gives the implementation of the generator. It
first attempts deterministic synthesis and falls back to LLM synthesis only when
deterministic synthesis returns no tree. Both paths use rendering and passing go build as
the acceptance gate---even if deterministic synthesis returns a tree, that is
not enough; it must actually compile, otherwise it is treated as unresolved and
passed to LLM synthesis. Both paths are constrained by the same set of
structural checks, and neither determines whether the code expresses the
normative rule, which is the job of the verification chain of
\S\ref{subsec:exec-verify}; they only confine the output to a traceable space.

\textbf{Deterministic main path.} When an IR predicate selects an atomic
template through the predicate-to-template map and its subject and constraint can fill that template's
parameter slots, the generator deterministically derives a DSL tree using
predefined rules. Here, deriving means constructing a DSL tree from the IR step by
step by fixed rules, with no LLM call throughout. Most lintable normative rules in our
study are covered by this deterministic path.

\textbf{Restricted LLM synthesis.} When deterministic synthesis returns no tree,
the generator switches to LLM synthesis. The model receives the normative-rule context,
the structured IR, the candidate atomic-template set, and a readable enumeration
and signature table of $\mathcal{V}$ and $\mathcal{A}$. It
outputs only one of two results, either a serialized DSL tree or an explicit
no-template abstention marker NT. The raw LLM
output is first validated by a parser, which resolves it into
a legal DSL tree $t \in \Tv$, a type error, or NT, and any
identifier outside $\mathcal{V}/\mathcal{A}$ is rejected here; a passing tree is
rendered by $\rho$, and then a deterministic post-processor
binds Description/Citation/Name and the metadata of
each normative source, with details in Appendix~\ref{app:tv}.

\subsection{Three-Layer Semantic Alignment Verification}
\label{subsec:exec-verify}

Whether the generated code expresses the normative rule is the hardest judgment
in the whole pipeline. This subsection checks it at three levels from coarse to
fine.

\textbf{Layer~A, description provenance.} Check whether the Description
can be traced back to the normative source.

\textbf{Layer~B, semantic equivalence of the code-behavior summary.} It turns
the cross-modal judgment of whether the code implements the
normative rule into the same-modal judgment of whether two natural-language
sentences are semantically equivalent. Following $\smech$ of \S\ref{subsec:mech}, with the
phrase dictionary in Appendix~\ref{app:sigma}, the code is translated into a
code\_summary, which is then judged against the Description.

\textbf{Layer~C, compilation and structure.} Verify that the code can be parsed
syntactically, that the function and metadata fields are complete, that
dependencies are correctly imported, and that the lint is registered.

The three layers form a Boolean acceptance gate by conjunction. Layer~B's
judgment is the core signal, and Layer~A's provenance serves as a
prefilter; failing any layer triggers the repair loop of Appendix~\ref{sec:saiv}.
Since $\rho$ and $\smech$ are deterministic and Layer~A anchors the Description
to source text, the LLM contributes only this one same-modal comparison, and an
IR error surfaces here as a mismatch.

\textbf{Certificate-level execution verification, execution-level evidence for
\CodeEqIR.} Layer~B's semantic-equivalence judgment is given by an LLM and cannot be fully
determinized; as a complement, we introduce a model-independent
deterministic verification to check whether the generated code executes the DSL
tree it corresponds to, written \CodeEqIR. Its idea borrows
from the test oracle in software testing~\cite{barr2015oracle}, using a set of samples with
known-correct answers to check the code under test. It proceeds in two steps.

Step 1, checking a single atomic template. For each atomic template we
construct a pair of controlled certificate samples, one that should
make the template's predicate true, and one that should make it false. The
samples are constructed by the atomic template's parameter types and are not
bound to any concrete normative rule. The rendered code is executed for real on the two
certificates and the verdicts are read back; the atomic template is accepted by
this check if and only if the results on both certificates match expectation.

Step 2, generalizing from atomic templates to the whole tree. For a
generated lint holding a DSL tree $t$, if $t$ is composed only of accepted
atomic templates and compiles as a whole, we can induct over the tree
structure. The behavior of the leaves is checked in Step~1, and the combinators
$\neg/\wedge/\vee$ evaluate by classical propositional logic and preserve that
behavior level by level, so the rendered code $\rho(t)$ executes the whole tree,
written $\mathrm{Code}\equiv t$; and because $t$ is deterministically reduced from
the IR, \CodeEqIR.

This verification is performed on controlled certificate samples; the result is
deterministic and does not vary with run count or model. Its boundary
must be stressed. \CodeEqIR{} only shows that the code follows the intermediate
representation IR, which is not the same as code--normative-rule
alignment; if the IR itself has already drifted from the normative rule, this
verification cannot detect it. The code--normative-rule judgment therefore still
rests with \CodeEqRule{}, the Layer~B semantic-equivalence judgment.
